% EPL master thesis covers template
%documentclass{masterthesis}
\documentclass{EPL-master-thesis-covers-EN}
\raggedbottom

% = Packages ================================================================================
\usepackage[export]{adjustbox}  % frame around pictures
\usepackage{xspace}
\usepackage{xcolor}
\usepackage{geometry}
\usepackage{url}
\usepackage{hyperref}
\usepackage{minted}     % code listings
\setminted{%
  fontsize=\small,
  frame=lines,
  breaklines=true,
  breakanywhere=true,
  breakautoindent=true,
  breaksymbol={\color{gray}\footnotesize$\hookrightarrow$},
  tabsize=4,
}
\usepackage{import}
\usepackage{standalone}
\usepackage{graphicx}   % images
\usepackage{pdfpages}   % import pdfs
\usepackage{pdflscape}
\usepackage{lscape}
\usepackage{todonotes}  % \todo{}
\usepackage{enumitem}   % customize lists
\usepackage{subcaption} % subfigures
\usepackage{colortbl}   % color cells of a table
\usepackage{threeparttable} % tables with footnotes,...
\usepackage{multirow}   % cells that span over multiple rows in tables
\usepackage{pifont}     % icons \ding{}
\usepackage{glossaries} % glossary
\usepackage[nospace,nocompress]{cite}   % Fix \nocite errors
\usepackage{comment}
\usepackage{epsfig}
\usepackage{array}
\usepackage{supertabular}
\usepackage{gensymb}    % ° symbol
\usepackage{textcomp}   % \textcent
\usepackage{amssymb}    % math R symbol
\usepackage{longtable} % allows tables to span multiple pages
\usepackage{booktabs}  % optional, nicer horizontal lines
\usepackage{amsmath}   % math environments and \mathbf
\usepackage{tikz}
\usetikzlibrary{shapes.geometric, arrows, positioning, shadows}
\usepackage{pgfplots}
\pgfplotsset{compat=1.18}

\frenchspacing
\setlength{\emergencystretch}{3em}
\hbadness=10000
\vbadness=10000
\hfuzz=1000pt
\vfuzz=50pt
\newcolumntype{L}[1]{>{\raggedright\arraybackslash}p{#1}}

\hypersetup{
    colorlinks,
    linkcolor={red!50!black},
    citecolor={blue!50!black},
    urlcolor={blue!80!black}
}

% = Colors ==================================================================================
\definecolor{Blue}{rgb}{0,0,1}
\definecolor{Red}{rgb}{1,0,0}
\definecolor{Magenta}{rgb}{1,0,1}
\definecolor{dark-gray}{gray}{0.2}
\definecolor{light-gray}{gray}{0.95}
\definecolor{formalshadeOrange}{HTML}{FFF8F4}
\definecolor{DarkOrange}{RGB}{237, 130, 49}

% = Commands ================================================================================
\newcommand{\jean}[1]{\textcolor{Blue}{#1}}
\newcommand{\arthur}[1]{\textcolor{Red}{#1}}
\newcommand{\alex}[1]{\textcolor{Magenta}{#1}}
\newcommand{\eg}{\textit{e.g.,}\xspace}
\newcommand{\ie}{\textit{i.e.,}\xspace}
\newcommand{\etal}{\textit{et al.}\xspace}
\newcommand{\vs}{\textit{vs.}\xspace}

% Keep headings with some following lines to avoid orphaned section titles at page bottoms
\usepackage{etoolbox}
\preto\section{\Needspace{10\baselineskip}}
\preto\subsection{\Needspace{8\baselineskip}}
\preto\subsubsection{\Needspace{6\baselineskip}}

% = COVER ===================================================================================
% Title of the thesis
\title{Automatic Recognition and Generation of SNOMED codes in Medical Documents}

% Subtitle - remove this line if not applicable
%\subtitle{Development and Applications}

% Name of the student author(s)
\author{Guillaume \textsc{Gillet}}

% Official title of the master degree (copy/paste from list below)
\degreetitle{Master [120] in Computer Science and Engineering}

% Name of the supervisor(s)
\supervisor{Coralie \textsc{Hemptinne}}
\supervisor{Jean \textsc{Vanderdonckt}}

% Name of the reader(s)
\readerone{Maxime \textsc{Griot}}
\readertwo{Demet \textsc{Yuksel}}			% remove if not applicable

% Academic year (update if necessary)
\years{2025--2026}
\usepackage{titlesec}
\usepackage{needspace}
%\titleformat{\subsubsection}[runin]{}{}{}{}[]

% = GLOSSARY ================================================================================
\makeglossaries
    \newglossaryentry{ann}{name=ANN, description={Artificial Neural Network}}
    \newglossaryentry{api}{name=API, description={Application Programming Interface}}
    \newglossaryentry{cnn}{name=CNN, description={Convolutional Neural Network}}
    \newglossaryentry{crf}{name=CRF, description={Conditional Random Fields}}
    \newglossaryentry{ehr}{name=EHR, description={Electronic Health Record}}
    \newglossaryentry{faiss}{name=FAISS, description={Facebook AI Similarity Search}}
    \newglossaryentry{hmm}{name=HMM, description={Hidden Markov Model}}
    \newglossaryentry{knn}{name=k-NN, description={k-Nearest Neighbor}}
    \newglossaryentry{llm}{name=LLM, description={Large Language Model}}
    \newglossaryentry{loc}{name=LOC, description={Lines Of Code}} % NOTE: not used as abbreviation in body text; retained for completeness
    \newglossaryentry{miou}{name=mIoU, description={Mean Intersection over Union}}
    \newglossaryentry{mrr}{name=MRR, description={Mean Reciprocal Rank}}
    \newglossaryentry{ndcg}{name=NDCG, description={Normalized Discounted Cumulative Gain}}
    \newglossaryentry{ner}{name=NER, description={Named Entity Recognition}}
    \newglossaryentry{nlp}{name=NLP, description={Natural Language Processing}}
    \newglossaryentry{rag}{name=RAG, description={Retrieval-Augmented Generation}}
    \newglossaryentry{rnn}{name=RNN, description={Recurrent Neural Network}}
    \newglossaryentry{rf2}{name=RF2, description={Release Format 2 (SNOMED CT distribution format)}}
    \newglossaryentry{sdk}{name=SDK, description={Software Development Kit}}
    \newglossaryentry{svm}{name=SVM, description={Support Vector Machine}}
    \newglossaryentry{ucl}{name=UCLouvain, description={Université Catholique de Louvain}} % NOTE: not used as abbreviation in body text; retained for completeness
    \newglossaryentry{ui}{name=UI, description={User Interface}}
    \newglossaryentry{postcoord}{name=Post-coordination, description={The combination of multiple SNOMED CT concepts to represent a complex clinical idea not expressible by a single concept}}
    \newglossaryentry{biencoder}{name=Bi-encoder, description={A neural architecture in which query and candidate are encoded independently by the same or similar encoders, enabling efficient approximate nearest-neighbour search}}
    \newglossaryentry{entitylinking}{name=Entity linking, description={The task of associating a text mention (span) with its corresponding entry in a structured knowledge base or ontology}}
    \newglossaryentry{spandetection}{name=Span detection, description={The subtask of identifying the character offsets of a named entity mention within a text, prior to concept assignment}}

% = DOCUMENT ================================================================================
\begin{document}
    % ====================== FRONT COVER =====================
    \hypersetup{pageanchor=false}
    \maketitle
    %\restoregeometry
    %\newpage\null\thispagestyle{empty}\newpage % ajoute une page vierge
    \clearpage\thispagestyle{empty}\mbox{}\clearpage

    % ======================= ABSTRACT =======================
    \newgeometry{top=3.5cm,bottom=3.5cm,left=4cm,right=4cm}
    \begin{abstract}
    \thispagestyle{plain}
    \pagenumbering{arabic}

Manual SNOMED CT coding of clinical documents burdens physicians with a time-consuming, error-prone task that existing tools handle poorly for non-English text. At the Cliniques universitaires Saint-Luc in Brussels, ophthalmologists face this challenge daily: their French-language consultation notes are inadequately supported by the Epic EHR's built-in coding assistance, motivating the need for an automated solution.

This thesis implements and rigorously compares four computational approaches to automatic SNOMED CT entity linking: a three-tier dictionary with fuzzy matching (M1), a SapBERT bi-encoder with FAISS retrieval (M2), an LLM-driven named entity recognition pipeline grounded by SNOMED vector search (M3), and a hybrid architecture combining LLM span detection with tiered dictionary lookup (M4).

On the SNOMED CT Entity Linking Challenge 1.2.1 benchmark (272 MIMIC-IV discharge summaries, 75{,}491 gold annotations), M1 achieves an mIoU of 0.3601, competitive with the challenge winner (0.420), while LLM-based methods score near zero (0.016--0.040). This apparent underperformance is explained by what this thesis identifies as the \emph{annotation density paradox}: the benchmark exhaustively labels every medical term (${\sim}340$ per note), while LLM methods selectively extract clinically meaningful entities (${\sim}35$ per note). When measured independently of extraction density, M4 achieves the highest MRR (0.68) and a Hits@5 of 0.73 across all methods.

The French clinical evaluation reveals a complete reversal of rankings. On six anonymised ophthalmological PDFs from Saint-Luc, LLM-based methods outperform dictionary approaches: M3 and M4 identify 43.8\% of physician-validated diagnoses versus 37.5\% for M1. A targeted prompt improvement explicitly instructing the model to capture systemic comorbidities raises M4 to 56.3\% (M4*), surpassing both GPT-3.5 and Claude Sonnet 4.6 without retrieval grounding, a 12.5-point gain requiring no model retraining. This advantage is confirmed on 36 structured French clinical cases (270 gold annotations), where M4 achieves the highest top-5 hit rate (42\%).

These results demonstrate that benchmark performance alone is an insufficient basis for clinical deployment decisions. For Saint-Luc, M4* offers the best trade-off between coverage, concept precision, and SNOMED validity guarantees, all without transmitting patient data to external servers.


    \end{abstract}


        % =================== ACKNOWLEDGMENTS ====================

    % =================== TABLE OF CONTENTS ====================
    \tableofcontents
    \newpage
    \hypersetup{pageanchor=true}

    \chapter{Introduction}
    \subimport{chapters/Table/}{Introduction.tex}

    \chapter{Related Work}
    \subimport{chapters/Table/}{RelatedWork.tex}

    \chapter{Methodology}
    \subimport{chapters/Table/}{Methodology.tex}

    \chapter{Implementation}
    \subimport{chapters/Table/}{Implementation.tex}

    \chapter{Results and Discussion}
    \subimport{chapters/Table/}{Results.tex}

    \chapter{Conclusion and Future Work}
    \subimport{chapters/Table/}{Conclusion.tex}

    \listoffigures
    \listoftables

    % ======================= GLOSSARY =======================
    \glsadd{ann}\glsadd{api}\glsadd{cnn}\glsadd{crf}\glsadd{ehr}\glsadd{faiss}\glsadd{hmm}\glsadd{knn}\glsadd{llm}\glsadd{loc}\glsadd{miou}\glsadd{mrr}\glsadd{ndcg}\glsadd{ner}\glsadd{nlp}\glsadd{rag}\glsadd{rnn}\glsadd{rf2}\glsadd{sdk}\glsadd{svm}\glsadd{ucl}\glsadd{ui}\glsadd{postcoord}\glsadd{biencoder}\glsadd{entitylinking}\glsadd{spandetection}
    \printglossary[nonumberlist]

    % ===================== BIBLIOGRAPHY =====================
    \nocite{*} % Print all references regardless of whether they were cited in the poster or not
    \bibliographystyle{plain} % Plain referencing style
    \bibliography{references} % Use the example bibliography file sample.bib

    % ===================== BIBLIOGRAPHY =====================
    \appendix

    % =================== APPENDIX: PROMPTS ==================
    \chapter{LLM Prompts Used in Methods 3 and 4}
    \label{app:prompts}

    This appendix reproduces the exact prompts passed to \texttt{llama3.1:8b} via Ollama
    in Methods~3, 4, and~4*. Each prompt listing is split into a \textit{system message}
    (static, sent once per API call) and a \textit{user message} (dynamic, assembled per
    note or per span batch). Prompts are extracted verbatim from
    \texttt{Link/method3\_llm\_rag.py} and \texttt{Link/method4\_hybrid.py}.

    \section{Method 3 NER Prompt}
    \label{sec:m3ner}

    Used in Stage~1 of Method~3 (\texttt{\_llm\_extract\_spans} in
    \texttt{method3\_llm\_rag.py}). Notes longer than 12\,000 characters are split into
    overlapping chunks (500-character overlap); each chunk is sent as a separate call with
    the same system message. Extracted spans are deduplicated by normalised span text before
    offset recovery.

    \smallskip\noindent\textit{Note:} After the MIMIC-IV benchmark evaluation was completed,
    the \texttt{SCOPE} paragraph from M4* (Section~\ref{sec:m4starner}) was merged into
    \texttt{\_NER\_SYSTEM} in the production code. The prompt below reflects the version
    used for the results reported in Chapter~5.

    \medskip
    \noindent\textbf{System message:}

    \begin{minted}[frame=lines, framesep=2mm, fontsize=\footnotesize,
                   breaklines=true, bgcolor=light-gray]{text}
You are a clinical NLP expert. Extract all medical entities from the following
clinical note that could be mapped to SNOMED CT concepts.

CRITICAL RULES:
- Copy the span EXACTLY as it appears in the note -- character for character.
- NEVER expand abbreviations. If the note says 'IOL', return 'IOL'. If it says
  'CA profonde', return 'CA profonde'. If it says 'neuropathie optique bilat',
  return 'neuropathie optique bilat'.
- NEVER paraphrase, translate, or complete partial words.
- Short abbreviations (IOL, DMEK, OCT, OD, OG, AV, PIO, HTA, HTO) are valid
  entities -- include them.

Semantic types: disorder, finding, procedure, body_structure, substance,
observable_entity

Return ONLY a JSON array. Each element: {"span": "...", "type": "..."}
No explanation, no markdown fencing -- raw JSON only.
    \end{minted}

    \medskip
    \noindent\textbf{User message:}

    \begin{minted}[frame=lines, framesep=2mm, fontsize=\footnotesize,
                   breaklines=true, bgcolor=light-gray]{text}
{note_text_or_chunk}
    \end{minted}

    \section{Method 4 NER Prompt (Original Scope)}
    \label{sec:m4ner}

    Method~4 calls \texttt{\_llm\_extract\_spans} imported directly from
    \texttt{method3\_llm\_rag.py} (see \texttt{method4\_hybrid.py}, line~156). The NER
    system message is therefore \textbf{identical to Section~\ref{sec:m3ner}} for the
    original-scope variant used in the MIMIC-IV benchmark evaluation. It is listed
    separately here to mark the experimental boundary between M4 and M4*, and to confirm
    that the only change introduced in M4* is the addition of the \texttt{SCOPE} paragraph
    documented in Section~\ref{sec:m4starner}.

    \medskip
    \noindent\textbf{System message:} identical to Section~\ref{sec:m3ner}.

    \medskip
    \noindent\textbf{User message:} identical to Section~\ref{sec:m3ner}.

    \section{Method 4* NER Prompt (Expanded Comorbidity Scope)}
    \label{sec:m4starner}

    M4* modifies the NER system message by inserting a \texttt{SCOPE} paragraph immediately
    after the opening task description. The \texttt{CRITICAL RULES} block, semantic type
    list, and output format are unchanged. The inserted paragraph is the sole source of the
    12.5 percentage point improvement in comorbidity recall observed on the French qualitative
    evaluation (Section~5.2.3).

    \medskip
    \noindent\textbf{System message} (addition relative to Section~\ref{sec:m3ner} shown
    between the dashed lines):

    \begin{minted}[frame=lines, framesep=2mm, fontsize=\footnotesize,
                   breaklines=true, bgcolor=light-gray]{text}
You are a clinical NLP expert. Extract all medical entities from the following
clinical note that could be mapped to SNOMED CT concepts.

- - - - - - - - - - - - - - - - - - - - - - - - - - - - - - - - - - - - - - -
SCOPE: Include ALL diagnoses and findings -- both primary ophthalmic findings
AND systemic comorbidities mentioned anywhere in the note (patient history,
background conditions, co-morbidities such as hypertension, diabetes, etc.).
- - - - - - - - - - - - - - - - - - - - - - - - - - - - - - - - - - - - - - -

CRITICAL RULES:
- Copy the span EXACTLY as it appears in the note -- character for character.
- NEVER expand abbreviations. If the note says 'IOL', return 'IOL'. If it says
  'CA profonde', return 'CA profonde'. If it says 'neuropathie optique bilat',
  return 'neuropathie optique bilat'.
- NEVER paraphrase, translate, or complete partial words.
- Short abbreviations (IOL, DMEK, OCT, OD, OG, AV, PIO, HTA, HTO) are valid
  entities -- include them.

Semantic types: disorder, finding, procedure, body_structure, substance,
observable_entity

Return ONLY a JSON array. Each element: {"span": "...", "type": "..."}
No explanation, no markdown fencing -- raw JSON only.
    \end{minted}

    \medskip
    \noindent\textbf{User message:} identical to Section~\ref{sec:m3ner}.

    \section{Method 4 Reranking Prompt}
    \label{sec:m4rerank}

    Used when LLM reranking is enabled in Method~4 (\texttt{\_llm\_rerank\_batch} in
    \texttt{method4\_hybrid.py}). Spans are sent in batches of up to 8 per LLM call,
    reducing API overhead by approximately $8\times$ relative to one-span-per-call
    reranking. Up to 20 FAISS candidates are presented per span. The model returns a JSON
    array whose ordering must match the input batch order; the pipeline handles partial
    failures (missing or malformed entries) by retaining the pre-reranking FAISS rank for
    that span.

    Note: as documented in Section~5.5.5, LLM reranking was \textbf{not enabled} during
    the MIMIC-IV benchmark evaluation reported in this thesis. This prompt is provided for
    completeness and for use in the planned follow-up evaluation.

    \medskip
    \noindent\textbf{System message:}

    \begin{minted}[frame=lines, framesep=2mm, fontsize=\footnotesize,
                   breaklines=true, bgcolor=light-gray]{text}
You are a clinical terminology expert. For each numbered span below, select the
best SNOMED CT concept from its candidate list. Use the shared clinical context.

Return ONLY a JSON array (one object per span, same order):
[{"span_id": 0, "selected_concept_id": "..."},
 {"span_id": 1, "selected_concept_id": "..."}, ...]
No explanation, no markdown -- raw JSON array only.
    \end{minted}

    \medskip
    \noindent\textbf{User message} (one block per span in the batch, separated by blank
    lines; up to 8 spans per call):

    \begin{minted}[frame=lines, framesep=2mm, fontsize=\footnotesize,
                   breaklines=true, bgcolor=light-gray]{text}
[0] Span: "<verbatim span text>"
    Context: <up to 200 characters of surrounding note text>
    Candidates:
      1. [<concept_id>] <preferred term / description>
      2. [<concept_id>] <preferred term / description>
      ...up to 20 candidates

[1] Span: "<verbatim span text>"
    Context: <up to 200 characters of surrounding note text>
    Candidates:
      1. [<concept_id>] <preferred term / description>
      ...
    \end{minted}

    % ====================== BACK COVER =======================
    \newpage\null\thispagestyle{empty}\newpage % ajoute une page vierge
    \backcoverpage % Back cover page


\end{document}
